
%%%%%%%%%%%%%%%%%%%%%%%%%%%%%%%%%%%%%%%%
%           main body
%%%%%%%%%%%%%%%%%%%%%%%%%%%%%%%%%%%%%%%%


%-----------------------------------
%	TITLE PAGE
%-----------------------------------

\title[Probability \& Statistics]{\LARGE \bfseries 第三部分\quad 概率与统计} % The short title appears at the bottom of every slide, the full title is only on the title page
\author[Law of Large Numbers] % Your institution as it will appear on the bottom of every slide, may be shorthand to save space
{\Large \bfseries \S 6. 大数定律}
% \author{Singular Value Decomposition} % Your name
% \date{\today} % Date, can be changed to a custom date
\date{}

%------------------------------------------------

\begin{document}


%------------------------------------------------

\begin{frame}

\titlepage % Print the title page as the first slide

\end{frame}

%------------------------------------------------

\begin{frame}
\frametitle{内容提要} % Table of contents slide, comment this block out to remove it
\tableofcontents % Throughout your presentation, if you choose to use \section{} and \subsection{} commands, these will automatically be printed on this slide as an overview of your presentation
\end{frame}

%-----------------------------------
%	PRESENTATION SLIDES
%-----------------------------------

\section{切比雪夫不等式}

%------------------------------------------------

\begin{frame}

\begin{block}{切比雪夫不等式，Chebyshev's Inequality}
设随机变量$X$的期望$E(X)$与方差$Var(X)$都存在，那么对任意的正实数$\varepsilon > 0$, 下列不等式都成立
$$P(\vert X - E(X) \vert \ge \varepsilon) \le \dfrac{Var(X)}{\varepsilon^2}$$
\end{block}

\end{frame}

%------------------------------------------------

\begin{frame}

\begin{block}{切比雪夫不等式离散型情况的证明}
设$X$是离散型随机变量，且有分布列$p_i = P(X = x_i)$, 那么
\begin{align*}
& P(\vert X - E(X) \vert \ge \varepsilon) \\
= & \sum\limits_{\vert x_i - E(X) \vert \ge \varepsilon} p_i \le \sum\limits_{\vert x_i - E(X) \vert \ge \varepsilon} p_i \cdot \dfrac{(x_i - E(X))^2}{\varepsilon^2} \\
= & \dfrac{1}{\varepsilon^2} \sum\limits_{\vert x_i - E(X) \vert \ge \varepsilon} p_i \cdot (x_i - E(X))^2 \\
\le & \dfrac{1}{\varepsilon^2} \sum\limits_i p_i \cdot (x_i - E(X))^2 = \dfrac{Var(X)}{\varepsilon^2}
\end{align*}
\end{block}

\pause

连续型随机变量的情形也能类似证明，大家自行推导。

\end{frame}

%------------------------------------------------

\section{大数定律}

%------------------------------------------------

\begin{frame}

\begin{block}{依概率收敛，Convergence in Probability}
设$X_1,\ldots,X_n,\ldots$为一列随机变量，$a\in \mathbb{R}$为一个常数，若有
$$\lim\limits_{n\to\infty} P(\vert X_n - a \vert < \varepsilon) = 1$$
对任意正实数$\varepsilon$成立，那么称随机变量序列$\{X_n\}_{n\ge1}$依概率收敛于$a$, 并记作
$$X_n \xrightarrow{P} a$$

\pause

自然，我们可以类似地定义随机变量序列依概率收敛于一个随机变量
$$X_n \xrightarrow{P} X$$
\end{block}

\end{frame}

%------------------------------------------------

\begin{frame}

\begin{block}{切比雪夫定理}
设互相独立的随机变量序列$\{X_n\}_{n\ge1}$有相同的期望与方差，即
$$E(X_n) = \mu, \; Var(X_n) = \sigma^2, \quad n=1,2,\ldots$$
那么$\dfrac{1}{n}\sum\limits_{i=1}^n X_i$依概率收敛于$\mu$, 也就是说
$$\lim\limits_{n\to\infty} P\left( \left\vert \dfrac{1}{n}\sum\limits_{i=1}^n X_i - \mu \right\vert < \varepsilon \right) = 1$$
对任意正实数$\varepsilon$成立。
\end{block}

\end{frame}

%------------------------------------------------

\begin{frame}

\begin{block}{伯努利大数定律，Law of Large Numbers, LLN}
设$n_A$是$n$次独立重复试验中事件$A$的出现次数，$p$为每次试验中事件$A$出现的概率（即参数$p$的伯努利试验）。那么对任意的正实数$\varepsilon$, 有
$$\lim\limits_{n\to\infty} P\left( \left\vert \dfrac{n_A}{n} - \mu \right\vert < \varepsilon \right) = 1$$
\end{block}

\vspace{1em}
\pause

伯努利大数定律是切比雪夫定理的一个特殊情况，也就是说令
$$X_n = \begin{cases}
1, & \text{若第$n$次试验中事件$A$发生} \\
0, & \text{若第$n$次试验中事件$A$不发生} 
\end{cases} \quad n=1,2,\ldots$$
那么随机变量序列$\{X_n\}_{n\ge1}$满足切比雪夫定理条件。
\end{frame}

%------------------------------------------------

\begin{frame}

\begin{block}{伯努利大数定律的意义}
\begin{itemize}
\item 当事件$A$发生的概率$p$已知时，大数定律可以用来预测试验结果。
\item 当事件$A$发生的概率$p$未知时，因为当试验次数充分大时，事件$A$发生的频率$\dfrac{n_A}{n}$与事件$A$发生的概率$p$差别很小（依概率收敛），于是可以用事件$A$发生的频率来估计概率。此即所谓的``逆概率问题''，是 Monte Carlo 方法的起源。
\item 如果事件$A$发生的概率很小，则由大数定律知道其发生的频率也是很低的，或者说事件$A$很少发生，以至于在一次试验中几乎不会发生。这就是所谓的``小概率原理''。
\end{itemize}
\end{block}

\vspace{1em}
\pause

伯努利大数定律的条件可以一定程度的减弱。能减弱到什么程度，依然处在研究中。目前各种形式的大数定律的研究还远未达到完善的程度。

\end{frame}

%------------------------------------------------

\begin{frame}

\begin{block}{更一般的大数定律}
设$\{X_n\}_{n\ge1}$是一列随机变量，如果下列条件之一成立
\begin{enumerate}
\item 切比雪夫条件：随机变量两两不相关，期望、方差都存在，且方差有界
\item 马尔科夫（Markov）条件：$\lim\limits_{n\to\infty} \dfrac{1}{n}Var\left( \sum\limits_{i=1}^n X_i \right) = 0$
\end{enumerate}
那么对于任意正实数$\varepsilon$, 有
$$\lim\limits_{n\to\infty} P\left( \left\vert \dfrac{1}{n}\sum\limits_{i=1}^n X_i - \dfrac{1}{n}\sum\limits_{i=1}^n E(X_i) \right\vert < \varepsilon \right) = 1$$
\end{block}

\end{frame}

%------------------------------------------------

\begin{frame}

\begin{block}{辛钦（Khinchin）大数定律}
设$\{X_n\}_{n\ge1}$是一列独立同分布的随机变量，$E(X_n) = \mu$.那么对任意正实数$\varepsilon$, 有
$$\lim\limits_{n\to\infty} P\left( \left\vert \dfrac{1}{n}\sum\limits_{i=1}^n X_i - \mu \right\vert < \varepsilon \right) = 1$$
\end{block}

\vspace{1em}
\pause

注意与切比雪夫定理的区别。这里对随机变量的方差并{\color{red}没有加任何假设}。

\end{frame}

%------------------------------------------------

% \begin{frame}

% \begin{block}{Vandermonde行列式的计算}
% \begin{enumerate}
% % \addtocounter{enumi}{0}
% \item 把第$n-1$行乘以$-a_1$加到第$n$行，再把第$n-2$行乘以$-a_1$加到第$n-1$行。按照这种顺序依次向上，直到用第$2$行乘以$-a_1$加到第$1$行。由行列式性质，$V_n$的值不变。
% \[
% V_n = \det \begin{pmatrix}
% 1 & 1 & \cdots & 1 \\
% 0 & a_2-a_1 & \cdots & a_n-a_1 \\
% \vdots & \vdots & & \vdots \\
% 0 & a_2^{n-2}(a_2-a_1) & \cdots & a_n^{n-2}(a_n-a_1) \\
% \end{pmatrix}
% \]
% \end{enumerate}
% \end{block}

% \end{frame}

%------------------------------------------------
% % closing page
% \begin{frame}
% \HUGE{\centerline{\bfseries {\color{gray} The} {\color{zkblue} End}}}

% \pause

% \vspace{1.2em}

% \Large{\centerline{\bfseries {\color{gray}Thank} \quad {\color{zkblue} You}}}

% \end{frame}

%----------------------------------------------------------------------------------------

\end{document}
